\documentclass[11pt,oneside,spanish]{article}

%%%%%%%%%%%%%%%%%%%%%%%%%%%%%%%%%%%%%%%%%%%%%%%%%%%%%%%

\usepackage[T1]{fontenc}
\usepackage[latin1]{inputenc}
\usepackage[spanish]{babel}
\usepackage{array}
\usepackage{float}
\usepackage{amsmath}
\usepackage{amsfonts}
\usepackage{amssymb}
\usepackage{amstext}
\usepackage{amsthm}
\usepackage{enumerate}
\usepackage{fancyhdr}
\usepackage[usenames,dvipsnames,svgnames,table]{xcolor}
\usepackage{tikz}
\usepackage[linkcolor=blue,colorlinks=true,urlcolor=blue]{hyperref}

\usepackage{import}
\graphicspath{{figuras/}}

\usepackage{multicol}

\usepackage{enumerate}

\usepackage{booktabs}

\theoremstyle{definition}
\newtheorem{defi}{Definici\'o—n}
\newtheorem{car}{Caracterizaci\'o—n}
\headheight 36pt
\usepackage{geometry}
\geometry{verbose,letterpaper,tmargin=30mm,bmargin=30mm,lmargin=25mm,rmargin=25mm}
\parindent 0em
\parskip 0ex

\newtheorem{ejercicio}{Ejercicio}
\newtheorem{teorema}{Teorema}
\newtheorem{lema}{Lema}
\newtheorem{corolario}{Corolario}
\theoremstyle{definition}\newtheorem{definicion}{Definici—n}
\theoremstyle{definition}\newtheorem{ejemplo}{Ejemplo}
\theoremstyle{remark}\newtheorem{nota}{\textsc{Nota}}
\theoremstyle{definition}\newtheorem{proposicion}{Proposici—n}
\theoremstyle{definition}\newtheorem{problema}{Problema}
\newenvironment{demostracion}{\textsc{Demostraci\'{o}n.}}{\halmos}

\newcommand{\bb}[1]{\mathbb{#1}}
\newcommand{\eq}[1]{\begin{equation} #1 \end{equation}}
\newcommand{\dpr}[2]{\frac{\partial #1}{\partial #2}}
\newcommand{\R}{\mathbb{R}}
\newcommand{\K}{\mathbb{K}}
\newcommand{\N}{\mathbb{N}}
\newcommand{\I}{\mathbb{I}}
\newcommand{\emqc}[1]{#1 \in \N}
\newcommand{\ssuc}{\{x_{n_k}\}_{k\in\N}}
\newcommand{\Z}{\mathbb{Z}}
\newcommand{\Q}{\mathbb{Q}}
\newcommand{\sen}{\text{sen}}
\newcommand{\senh}{\text{senh}}
\newcommand{\adh}[1]{\text{adh}({#1})}
\newcommand{\der}[1]{\text{der}({#1})}
\newcommand{\inte}[1]{\text{int}({#1})}
\newcommand{\fr}[1]{\text{fr}({#1})}
\newcommand{\co}[1]{\text{co}({#1})}
\newenvironment{respuesta}{\smallskip \textbf{\textsf{Respuesta}} \sffamily \newline}{\hfill}
\newcommand{\re}[1]{\begin{respuesta} #1 \end{respuesta}}  
\newenvironment{solucion}{\smallskip \textbf{\textsf{Soluci\'on}} \sffamily \newline}{\hfill}
\newcommand{\so}[1]{\begin{solucion} #1 \end{solucion}}  

\makeatletter
\def\ScaleIfNeeded{
\ifdim\Gin@nat@width>\linewidth
\linewidth
\else
\Gin@nat@width
\fi
}
\makeatother

\usepackage[mathlf,textlf,minionint,openg]{MinionPro}
\usepackage[bb=lucida,bbscaled=.95]{mathalfa}
\usepackage[protrusion=true,expansion=true]{microtype}
\renewcommand{\vec}[1]{\boldsymbol{#1}}

\usepackage{listings}

%%%%%%%%%%%%%%%%%%%%%%%%%%%%%%%%%%%%%%%%%%%%%%%%%%%%%%%

\begin{document}

\thispagestyle{fancy}
\lhead{Universidad Nacional Andr\'es Bello}
\rhead{An\'alisis de Datos}

\includegraphics[width=8em]{by-nc-sa.pdf}

\begin{center}
\textbf{Control de Laboratorio N$^\circ$1  - Secci\'on 1} 
\end{center}

\bigskip

Nombre: \underline{\textcolor{white}{---------------------------------------------------}} \hspace{2mm} RUN: \underline{\textcolor{white}{---------------------------}}

\bigskip

Escriba el c\'odigo o el desarrollo algebraico que permite resolver cada pregunta. Las preguntas sin justificaci\'on se consideran omitidas. Marque sus respuestas con \textbf{l\'apiz pasta o tinta} a continuaci\'on.

\bigskip

\begin{tikzpicture}[font=\small]
    \foreach \line in {1,2,...,10} {
        \begin{scope}[yshift=-\line cm]
            \foreach \count/\desc in {1/A, 2/B, 3/C, 4/D} { 
                \node at ({\count * 0.5}, 0.4) {\desc};
                \node at (0,0) {\normalsize\textbf{\line}};
                \node[draw,circle,inner sep=1pt] at ({\count * 0.5},0) {\count};
            }
        \end{scope}
    }
\end{tikzpicture}

\newpage

\begin{enumerate}[(1)]
\item Considere el siguiente vector de datos y el vector de ponderadores
\begin{lstlisting}[backgroundcolor=\color{Gray!20},frame=none,basicstyle=\ttfamily]
 x <- c(0.18, -1.54, 0.42, 0.95)
 w <- c(2, 1, 3, 1)
\end{lstlisting}
{\textquestiondown}Cu\'al es el valor de $\mu$ que minimiza la ecuaci\'on de m\'inimos cuadrados $\sum_{i=1}^n w_i(x_i-\mu)^2$?
	\begin{enumerate}[(a)]
	\item 0,1471
	\item 1,077
	\item 0,3
	\item 0,0025
	\end{enumerate}
	
\item Considere los siguientes datos
\begin{lstlisting}[backgroundcolor=\color{Gray!20},frame=none,basicstyle=\ttfamily]
 x <- c(0.8, 0.47, 0.51, 0.73, 0.36, 0.58, 0.57, 0.85, 0.44)
 y <- c(1.39, 0.72, 1.55, 0.48, -1.59, 1.23, -0.65, 1.49, 0.05)
\end{lstlisting}
Estime una regresi\'on que pase por el origen considerando $x$ como input e $y$ como output. (Hint: no centre los datos ya que se pide una regresi\'on 	que pase por el origen y no por las medias de los datos.)
{\textquestiondown}Cu\'al es el valor de la pendiente?
	\begin{enumerate}[(a)]
	\item 0,8623
	\item 1,0861
	\item -1,713
	\item 0,59915
	\end{enumerate}
	
\item Cargue la  base da datos \texttt{mtcars} y estime una regresi\'on que explique la variable \texttt{mpg} en funci\'on de la variable \texttt{weight}. {\textquestiondown}Cu\'al es el coeficiente de la pendiente?	
	\begin{enumerate}[(a)]
	\item -9,559
	\item 30,2851
	\item -5,344
	\item 0,5591
	\end{enumerate}

\item Considere una regresi\'on hipot\'etica en la cual $Y$ es el output y $X$ es el input. La desviaci\'on t\'ipica del input es la mitad de la del output. La correlaci\'on entre ambas variables es $0,5$. {\textquestiondown}Cu\'al es el valor de la pendiente de un modelo en el que se invierte el input y el output?
	\begin{enumerate}[(a)]
	\item 2
	\item 0,25
	\item 4
	\item 1
	\end{enumerate}

\newpage

\item Considere el vector 
\begin{lstlisting}[backgroundcolor=\color{Gray!20},frame=none,basicstyle=\ttfamily]
 x <- c(8.58, 10.46, 9.01, 9.64, 8.86)
\end{lstlisting}
Si se normalizan los datos {\textquestiondown}Cu\'al es el valor de la primera observaci\'on?	
	\begin{enumerate}[(a)]
	\item -0,9719
	\item 8,58
	\item 8,886
	\item 8,931
	\end{enumerate}

\item Sea $\beta_1$ la pendiente de un modelo en el que $Y$ es el output y $X$ es el input. Sea $\gamma_1$ la pendiente de un modelo en el que $X$ es el output e $Y$ es el input. {\textquestiondown}Cu\'al es el valor de $\beta_1 / \gamma_1$?		
	\begin{enumerate}[(a)]
	\item $\text{sd}(X)/\text{sd}(Y)$
	\item $\text{var}(Y)/\text{var}(X)$
	\item $1$
	\item $\text{cor}(X,Y)$
	\end{enumerate}
	
\item Considere un modelo  lineal con datos hipot\'eticos en el que el input y el output tienen media 0. {\textquestiondown}Qu\'e se puede afirmar respecto del intercepto si se estima una regresi\'on?		
	\begin{enumerate}[(a)]
	\item No hay informaci\'on suficiente
	\item Tiene valor uno
	\item Debe ser id\'entico a 0
	\item Se indefine ya que se divide por cero
	\end{enumerate}	

\item Considere el vector
\begin{lstlisting}[backgroundcolor=\color{Gray!20},frame=none,basicstyle=\ttfamily]
 x <- c(0.8, 0.47, 0.51, 0.73, 0.36, 0.58, 0.57, 0.85, 0.44, 0.42)
\end{lstlisting}
{\textquestiondown}Cu\'al es el valor que minimiza la suma de las distancias el cuadrado entre estos puntos y el valor mismo?
	\begin{enumerate}[(a)]
	\item 0,571
	\item 0,8
	\item 0,36
	\item 0,573
	\end{enumerate}

\newpage

\item Considere los siguientes vectores con $x$ como input e $y$ como output
\begin{lstlisting}[backgroundcolor=\color{Gray!20},frame=none,basicstyle=\ttfamily]
 x <- c(0.61, 0.93, 0.83, 0.35, 0.54, 0.16, 0.91, 0.62, 0.62)
 y <- c(0.67, 0.84, 0.6, 0.18, 0.85, 0.47, 1.1, 0.65, 0.36)
\end{lstlisting}
{\textquestiondown}Cu\'al es el valor $p$ de test de dos colas para determinar si acaso el $\beta_1$ de la regresi\'on lineal es 0 o no?
	\begin{enumerate}[(a)]
	\item 0,05296
	\item 0,391
	\item 0,364
	\item 0,0243
	\end{enumerate}
	
\item {\textquestiondown}Cu\'al es la desviaci\'on t\'ipica residual en el problema anterior?
	\begin{enumerate}[(a)]
	\item 0,438
	\item 0,223
	\item 0,435
	\item 0,355
	\end{enumerate}	

\end{enumerate}

\bigskip

\textbf{Tiempo: 20 minutos}

\medskip

\textbf{0,6 puntos por cada respuesta correcta}\\
\textbf{-0,6 puntos por cada respuesta incorrecta}\\
\textbf{-1 punto por cada respuesta omitida}

\end{document}